% !TeX program = pdflatex
% Gallery starter; see TEMPLATE-README.txt for attribution and licenses.
\documentclass[11pt,a4paper]{article}
\usepackage[T1]{fontenc}
\usepackage{lmodern,microtype}
\usepackage[margin=24mm]{geometry}
\usepackage{graphicx,booktabs,tabularx,enumitem,amsmath,amssymb,xcolor,titlesec,fancyhdr}
\usepackage[hidelinks]{hyperref}
\definecolor{accent}{HTML}{535483}\definecolor{quiet}{HTML}{526171}
\setlength{\parindent}{0pt}\setlength{\parskip}{5pt}
\setlist[itemize]{leftmargin=*,itemsep=3pt}
\titleformat{\section}{\Large\sffamily\bfseries\color{accent}}{\thesection}{0.6em}{}
\titleformat{\subsection}{\sffamily\bfseries}{\thesubsection}{0.6em}{}
\titlespacing*{\section}{0pt}{12pt}{5pt}
\titlespacing*{\subsection}{0pt}{10pt}{4pt}
\pagestyle{fancy}\fancyhf{}\fancyhead[L]{\small\sffamily Theorem and Proof Notes}\fancyhead[R]{\small\sffamily Alex Morgan}\fancyfoot[R]{\small\thepage}
\renewcommand{\headrulewidth}{0pt}
\newcommand{\eyebrow}[1]{{\small\sffamily\bfseries\color{accent}#1}\par\vspace{5mm}}
\newcommand{\hero}[2]{{\fontsize{32}{36}\selectfont\sffamily\bfseries #1\par}\vspace{4mm}{\large\color{quiet}#2\par}\vspace{6mm}{\color{accent}\rule{\linewidth}{1.5pt}}\vspace{5mm}}
\usepackage{amsthm}\newtheorem{theorem}{Theorem}[section]\newtheorem{definition}[theorem]{Definition}
\begin{document}
\eyebrow{BOOKS \& NOTES / 2026}
\hero{The geometry of projection}{Mathematical notes / Inner product spaces}
\section{Definitions}
\begin{definition}For a subspace $U$ of an inner product space $V$, a vector $v$ is orthogonal to $U$ if $\langle v,u\rangle=0$ for every $u\in U$.\end{definition}
\begin{theorem}[Orthogonal decomposition]Let $U$ be a finite-dimensional subspace of an inner product space $V$. Every $v\in V$ has a unique decomposition $v=u+w$, where $u\in U$ and $w\in U^\perp$.\end{theorem}
\begin{proof}Choose an orthonormal basis $e_1,\ldots,e_k$ of $U$ and set
\[
u=\sum_{j=1}^{k}\langle v,e_j\rangle e_j,\qquad w=v-u.
\]
Then $w$ is orthogonal to each basis vector. For uniqueness, subtract two decompositions and use $U\cap U^\perp=\{0\}$.\end{proof}
\section{An example}Project $(2,1)$ onto the line generated by $(1,1)$. The projection is $(3/2,3/2)$ and the residual $(1/2,-1/2)$ is orthogonal to the line.
\section{Exercises}\begin{enumerate}\item Verify the residual calculation.\item Prove that an orthogonal projection is idempotent.\item Explain the connection to least squares.\end{enumerate}
\end{document}
